\documentclass{article}

% -----------------------------
% Packages
% -----------------------------
\usepackage{fancyhdr}
\usepackage{extramarks}
\usepackage{amsmath}
\usepackage{amsthm}
\usepackage{amsfonts}
\usepackage{tikz}
\usepackage[plain]{algorithm}
\usepackage{algpseudocode}
\usepackage{amssymb}

\usetikzlibrary{automata,positioning}

% -----------------------------
% Basic Document Settings
% -----------------------------
\topmargin=-0.45in
\evensidemargin=0in
\oddsidemargin=0in
\textwidth=6.5in
\textheight=9.0in
\headsep=0.25in

\linespread{1.1}

\pagestyle{fancy}
\lhead{\hmwkAuthorName}
\chead{\hmwkClass: \hmwkTitle}
\rhead{}
\lfoot{\lastxmark}
\cfoot{\thepage}

\renewcommand\headrulewidth{0.4pt}
\renewcommand\footrulewidth{0.4pt}

\setlength\parindent{0pt}

% -----------------------------
% Problem Section Helpers
% -----------------------------
\newcommand{\enterProblemHeader}[1]{
  \nobreak\extramarks{}{Problem #1 continued on next page\ldots}\nobreak{}
  \nobreak\extramarks{Problem #1 (continued)}{Problem #1 continued on next page\ldots}\nobreak{}
}

\newcommand{\exitProblemHeader}[1]{
  \nobreak\extramarks{Problem #1 (continued)}{Problem #1 continued on next page\ldots}\nobreak{}
  \nobreak\extramarks{Problem #1}{}\nobreak{}
}

\setcounter{secnumdepth}{0}
\newcounter{partCounter}
\newcounter{homeworkProblemCounter}
\setcounter{homeworkProblemCounter}{1}
\nobreak\extramarks{Problem \arabic{homeworkProblemCounter}}{}\nobreak{}

% -----------------------------
% Theorem env (optional)
% -----------------------------
\newtheorem*{theorem}{Theorem}

% -----------------------------
% Homework Problem Environment
%   Usage:
%   \begin{homeworkProblem}            % auto-number
%     % Your answer here
%   \end{homeworkProblem}
%
%   \begin{homeworkProblem}[1.2.3]     % label problem "1.2.3" (e.g. textbook numbering)
%     % Your answer here
%   \end{homeworkProblem}
% -----------------------------
\newenvironment{homeworkProblem}[1][]{
  \def\hwProblemLabel{#1}
  \ifx\hwProblemLabel\empty
    \edef\hwProblemLabel{\arabic{homeworkProblemCounter}}
    \stepcounter{homeworkProblemCounter}
  \fi
  \section{Problem \hwProblemLabel}
  \setcounter{partCounter}{1}
  \enterProblemHeader{\hwProblemLabel}
}{
  \exitProblemHeader{\hwProblemLabel}
}

% -----------------------------
% Homework Details (EDIT THESE)
% -----------------------------
\newcommand{\hmwkTitle}{Homework \#6}
\newcommand{\hmwkClass}{Math 3320}
\newcommand{\hmwkAuthorName}{\textbf{Joseph Quinn}}

% -----------------------------
% Title
% -----------------------------
\title{
  \vspace{2in}
  \textmd{\textbf{\hmwkClass:\ \hmwkTitle}}\\
  \vspace{3in}
}
\author{\hmwkAuthorName}
\date{}

% -----------------------------
% Convenience Macros (optional)
% -----------------------------
\renewcommand{\part}[1]{\textbf{\large Part \Alph{partCounter}}\stepcounter{partCounter}\\}

% Algorithms
\newcommand{\alg}[1]{\textsc{\bfseries \footnotesize #1}}

% Calculus
\newcommand{\deriv}[1]{\frac{\mathrm{d}}{\mathrm{d}x} (#1)}
\newcommand{\pderiv}[2]{\frac{\partial}{\partial #1} (#2)}
\newcommand{\dx}{\mathrm{d}x}

% Statistics
\newcommand{\E}{\mathrm{E}}
\newcommand{\Var}{\mathrm{Var}}
\newcommand{\Cov}{\mathrm{Cov}}
\newcommand{\Bias}{\mathrm{Bias}}

% Proof step (for aligned derivations)
\newcommand{\step}[2]{& #1 & & \text{#2} \\}

% Solution header (optional)
\newcommand{\solution}{\textbf{\large Solution}}

% -----------------------------
% Document
% -----------------------------
\begin{document}

\begin{homeworkProblem}[A12]
	The primal LP we are given is
	\[
		\begin{aligned}
			\max \quad & z = 14x_1+12x_2+13x_3+4x_4 \\
			\text{s.t.}\quad
			           & x_1+4x_2+x_3 \le 17        \\
			           & 3x_1+2x_3+x_4 \le 10       \\
			           & 6x_1+x_2+4x_3+5x_4 \le 27  \\
			           & x_1,x_2,x_3,x_4 \ge 0.
		\end{aligned}
	\]

	Becuase we are dealing with a mazimation primal problem we can formulate the variable signs for the dual variables like so,
	\begin{align*}
		(P1) \implies y_{1} \ge 0 \\
		(P2) \implies y_{2} \ge 0 \\
		(P3) \implies y_{3} \ge 0 \\
	\end{align*}


	Our $w$ optimization function becomes the linear combination of the RHS constants, being
	\[
		\quad w = 17y_1+10y_2+27y_3
	\]

	Then by transposing the constraint coefficnets we can formulate the full dual LP like so,
	\[
		\begin{aligned}
			\min \quad & w = 17y_1+10y_2+27y_3                  \\
			\text{s.t.}\quad
			           & y_1+3y_2+6y_3 \ge 14  &  & \text{(D1)} \\
			           & 4y_1+y_3 \ge 12       &  & \text{(D2)} \\
			           & y_1+2y_2+4y_3 \ge 13  &  & \text{(D3)} \\
			           & y_2+5y_3 \ge 4        &  & \text{(D4)} \\
			           & y_1,y_2,y_3 \ge 0.
		\end{aligned}
	\]


	The optimal primal dictionary we are given is so,
	\begin{align*}
		x_2 & = 3 + \frac{1}{8}x_1 + \frac{1}{8}x_4 - \frac{1}{4}x_5 + \frac{1}{8}x_6   \\
		x_3 & = 5 - \frac{3}{2}x_1 - \frac{1}{2}x_4 - \frac{1}{2}x_6                    \\
		x_7 & = 4 - \frac{1}{8}x_1 - \frac{25}{8}x_4 + \frac{1}{4}x_5 + \frac{15}{8}x_6 \\
		z   & = 101 - 4x_1 - x_4 - 3x_5 - 5x_6
	\end{align*}

	So the optimal $x$ is the basic feasible solution from this dictionary which is,
	\[
		(x_{1}=0,x_{2}=3,x_{3}=5,x_{4}=0)
	\]

	Resulting in a $z$ value of 101

	Through the correspondence of dictionaries we know that each basic variable of the dual corresponds to a non basic variable from the primal, and a non basic dual variable to a basic primal variable.
	And through the formulation of the dual we know the coefficent attached to each $z$ equation $x$ variable becomes its corresponding $y$ dual variables value negated in the dual dictionary.

	we can write hte vairable correspondence for this LP like so,
	\begin{align*}
		y_{1} = x_{5} \\
		y_{2}=x_{6}   \\
		y_{3}=x_{7}   \\
		x_{1}=y_{4}   \\
		x_{2}=y_{5}   \\
		x_{3}=y_{6}   \\
		x_{4}=y_{7}
	\end{align*}

	Then from the prime optimal dicitonary we know $x_{2}, x_{3},x_{7}$ are basic, so their corresponding dual variables $y_{5},y_{6}, y_{3}$ will be non basic and have a value of 0.

	And on the other side the primal non basic $x_{1},x_{4},x_{5},x_{6}$ will become basic in the dual with the corresponding dual vairables $y_{4},y_{7},y_{1},y_{2}$. The values of which we stated earlier are the
	negated coefficents so the values become $y_{4}= 4,y_{7}=1,y_{1}=3,y_{2}=5$. \\

	So we can formulate $y*$ as $(y_{1}=3,y_{2}=5,y_{3}=0)$ \\

	Now to verify all dual constrints by substituting $y_1=3,\ y_2=5,\ y_3=0$ into every original dual constraint. \\

	Constraint (D1): $y_1+3y_2+6y_3 = 3+3(5)+6(0) = 18 \ge 14$

	Constraint (D2): $4y_1+y_3 = 4(3)+0 = 12 \ge 12$

	Constraint (D3):$y_1+2y_2+4y_3 = 3+2(5)+4(0) = 13 \ge 13$

	Constraint (D4): $y_2+5y_3 = 5+5(0) = 5 \ge 4 $ \\

	And checking nonnegativity constraints:
	\[
		y_1 = 3 \ge 0 , \qquad y_2 = 5 \ge 0 , \qquad y_3 = 0 \ge 0 .
	\]
	\\
	Therefore, $y^* = (3,5,0)$ is dual feasible. \\

	Finally lets verify the optimal value equality between the primal and the dual, the optimal primal dictionary gives $z=101$. We will substitute $y^* = (3,5,0)$ into the original dual objective function:
	\[
		w(y^*) = 17y_1+10y_2+27y_3 = 17(3)+10(5)+27(0) = 51+50+0 = 101.
	\]

	Thus,
	\[
		z(x^*) = w(y^*) = 101.
	\]

\end{homeworkProblem}

\begin{homeworkProblem}[A13]
	Given the LP problem we want to see if the following $x$ gives an optimal solution
	\[
		\begin{aligned}
			\max \quad z ={} &
			1.62x_{\mathrm{AA}} + 1.62x_{\mathrm{CA}}
			+ 1.62x_{\mathrm{SA}} + 1.62x_{\mathrm{IA}}                                         \\
			                 & + 1.08x_{\mathrm{AB}} + 1.08x_{\mathrm{CB}}
			+ 1.08x_{\mathrm{SB}} + 1.08x_{\mathrm{IB}}                                         \\[4pt]
			\text{s.t.}\quad
			                 & x_{\mathrm{AA}} + x_{\mathrm{AB}} \le 3814                       \\
			                 & x_{\mathrm{CA}} + x_{\mathrm{CB}} \le 2666                       \\
			                 & x_{\mathrm{SA}} + x_{\mathrm{SB}} \le 4016                       \\
			                 & x_{\mathrm{IA}} + x_{\mathrm{IB}} \le 1300                       \\
			                 & -7x_{\mathrm{AA}} + 7x_{\mathrm{CA}}
			+ 13x_{\mathrm{SA}} - 8x_{\mathrm{IA}} \le 0                                        \\
			                 & -16x_{\mathrm{AB}} - 2x_{\mathrm{CB}}
			+ 4x_{\mathrm{SB}} - 17x_{\mathrm{IB}} \le 0                                        \\
			                 & -2x_{\mathrm{AA}} + x_{\mathrm{CA}}
			- 3x_{\mathrm{SA}} + 14x_{\mathrm{IA}} \le 0                                        \\
			                 & -2x_{\mathrm{AB}} + x_{\mathrm{CB}}
			- 3x_{\mathrm{SB}} + 14x_{\mathrm{IB}} \le 0                                        \\[4pt]
			                 & x_{\mathrm{AA}},x_{\mathrm{CA}},x_{\mathrm{SA}},x_{\mathrm{IA}},
			x_{\mathrm{AB}},x_{\mathrm{CB}},x_{\mathrm{SB}},x_{\mathrm{IB}}
			\ge 0.
		\end{aligned}
	\]

	\begin{align*}
		X_AA = 3754 \\
		X_CA = 2666 \\
		X_SA = 920  \\
		X_IA = 543  \\
		X_AB = 60   \\
		X_CB = 0    \\
		X_SB = 3096 \\
		X_IB = 672
	\end{align*}

	First lets check (i) of CSC to validate that our $x$ is primal feasible \\

	We see all values are $\ge 0$ so it passes non negativity

	\begin{align*}
		X_AA = 3754 \ge 0 \\
		X_CA = 2666 \ge 0 \\
		X_SA = 920  \ge 0 \\
		X_IA = 543  \ge 0 \\
		X_AB = 60   \ge 0 \\
		X_CB = 0    \ge 0 \\
		X_SB = 3096 \ge 0 \\
		X_IB = 672 \ge 0
	\end{align*}

	And substituting $x$ into the LP's constraints we get:

	\begin{align*}
		(P1):\quad & 3754 + 60 = 3814 \le 3814                 \\
		(P2):\quad & 2666 + 0 = 2666 \le 2666                  \\
		(P3):\quad & 920 + 3096 = 4016 \le 4016                \\
		(P4):\quad & 543 + 672 = 1215 \le 1300                 \\
		(P5):\quad & -7(3754)+7(2666)+13(920)-8(543) = 0 \le 0 \\
		(P6):\quad & -16(60)-2(0)+4(3096)-17(672) = 0 \le 0    \\
		(P7):\quad & -2(3754)+2666-3(920)+14(543) = 0 \le 0    \\
		(P8):\quad & -2(60)+0-3(3096)+14(672) = 0 \le 0
	\end{align*}

	All primal constraints are satisfied, and all variables
	are nonnegative. Therefore, \(x\) is primal feasible. \\

	Now we can move to (ii) of CSC, looking at the above calculations after substituting in $x$ we see that
	only $P4$ is slack, this tells us the associated $y_4$ dual variable has to be tight and has a value of 0. We dont get any new information from the tight constraints.

	This helps us simpify the next step being (iv) of CSC\\

	The first step here is to form the Dual of the primal LP, since we are dealing with a standard form LP we know all hte constarints will have the $\ge$ signage and the
	variable constarints will all be $\ge 0$

	The optimzation of the dual will be minimization and the function will come from the RHS constants making it  $w  = 3814y_1 + 2666y_2 + 4016y_3 + 1300y_4 $ \\

	Then we transpose and negate the rest of the coefficients resulting in the Dual formation of
	\begin{align*}
		\min \quad w & = 3814y_1 + 2666y_2 + 4016y_3 + 1300y_4 \\
		\text{s.t.} \quad
		             & y_1 - 7y_5 - 2y_7 \ge 1.62              \\
		             & y_2 + 7y_5 + y_7 \ge 1.62               \\
		             & y_3 + 13y_5 - 3y_7 \ge 1.62             \\
		             & y_4 - 8y_5 + 14y_7 \ge 1.62             \\
		             & y_1 - 16y_6 - 2y_8 \ge 1.08             \\
		             & y_2 - 2y_6 + y_8 \ge 1.08               \\
		             & y_3 + 4y_6 - 3y_8 \ge 1.08              \\
		             & y_4 - 17y_6 + 14y_8 \ge 1.08            \\
		             & y_1,y_2,y_3,y_4,y_5,y_6,y_7,y_8 \ge 0
	\end{align*}

	Now continuing with the CSC we know all the $x$ variables besides $X_CB$ are slack, meaning there corresponding Dual equation must be tight, this gives us the information of,

	\begin{align*}
		(D1):\quad & y_1 - 7y_5 - 2y_7 = 1.62   \\
		(D2):\quad & y_2 + 7y_5 + y_7 = 1.62    \\
		(D3):\quad & y_3 + 13y_5 - 3y_7 = 1.62  \\
		(D4):\quad & y_4 - 8y_5 + 14y_7 = 1.62  \\
		(D5):\quad & y_1 - 16y_6 - 2y_8 = 1.08  \\
		(D7):\quad & y_3 + 4y_6 - 3y_8 = 1.08   \\
		(D8):\quad & y_4 - 17y_6 + 14y_8 = 1.08
	\end{align*}

	and we know that \(y_4=0\),

	Now we can solve hte above system of equations via mathematica resulting in
	\[
		\begin{aligned}
			y_1 & = 2.34, & y_5 & = 0.06, \\
			y_2 & = 1.05, & y_6 & = 0.06, \\
			y_3 & = 1.29, & y_7 & = 0.15, \\
			y_4 & = 0,    & y_8 & = 0.15.
		\end{aligned}
	\]


	This gives us our $y$ (potential) solution. We now have to check (ii) CSC to make sure its dual feasible,

	We see all values are $\ge 0$ so it passes nonnegativity.

	\[
		\begin{aligned}
			y_1 & = 2.34 \ge 0 \\
			y_2 & = 1.05 \ge 0 \\
			y_3 & = 1.29 \ge 0 \\
			y_4 & = 0 \ge 0    \\
			y_5 & = 0.06 \ge 0 \\
			y_6 & = 0.06 \ge 0 \\
			y_7 & = 0.15 \ge 0 \\
			y_8 & = 0.15 \ge 0
		\end{aligned}
	\]

	And substituting $y$ into the dual LP's constraints we get:

	\begin{align*}
		(D1):\quad & 2.34-7(0.06)-2(0.15)=1.62\ge1.62  \\
		(D2):\quad & 1.05+7(0.06)+0.15=1.62\ge1.62     \\
		(D3):\quad & 1.29+13(0.06)-3(0.15)=1.62\ge1.62 \\
		(D4):\quad & 0-8(0.06)+14(0.15)=1.62\ge1.62    \\
		(D5):\quad & 2.34-16(0.06)-2(0.15)=1.08\ge1.08 \\
		(D6):\quad & 1.05-2(0.06)+0.15=1.08\ge1.08     \\
		(D7):\quad & 1.29+4(0.06)-3(0.15)=1.08\ge1.08  \\
		(D8):\quad & 0-17(0.06)+14(0.15)=1.08\ge1.08
	\end{align*}

	All dual constraints are satisfied, and all variables
	are nonnegative. Therefore, \(y\) is dual feasible.


	Thus $x$ passes all CSC properties and we know it is an optimal solution, we can do a final check to make sure both $z(x)$ and $w(y)$ equations are equal,



	\begin{align*}
		z(x) & = 1.62(3754+2666+920+543)   \\
		     & \quad + 1.08(60+0+3096+672) \\
		     & = 1.62(7883)+1.08(3828)     \\
		     & = 16904.70
	\end{align*}


	\begin{align*}
		w(y) & = 3814(2.34)+2666(1.05)   \\
		     & \quad +4016(1.29)+1300(0) \\
		     & = 16904.70
	\end{align*}

	$z(x)=w(y)=16904.70$

	And it is an optimal solution!

\end{homeworkProblem}

\end{document}


